\section{Direction periods and collinear affine energy}
\label{sec:affine-collinear-period-energy}

The common-kernel selection in
Section~\ref{sec:affine-common-kernel-triple-selection} leaves a collinear
concentration problem.  We now extract an exact arithmetic period along each
nonnegative ray.  Write
\[
 U(h,k)=1+Rh,\quad V(h,k)=1+R(h+Ck),
 \quad W(h,k)=1+R(h+Bk)
\]
and parametrize the ray by
$x(n)=(h_0+ns,k_0+nt)$.  For a positive pairwise-coprime label
$\lambda=(d_U,d_V,d_W)$ define
\begin{align*}
 A_U&=s,&A_V&=s+Ct,&A_W&=s+Bt,\\
 r_Z&=\frac{d_Z}{\gcd(d_Z,A_Z)},&
 T_\lambda&=r_Ur_Vr_W,\\
 C_\lambda&=\gcd(d_U,A_U)\gcd(d_V,A_V)\gcd(d_W,A_W).
\end{align*}

\begin{theorem}[Exact ray period and capacity]
\label{thm:affine-exact-ray-period}
If the same label divides the corresponding three arms at $x(n)$ and
$x(m)$, then
\[
 T_\lambda\mid |n-m|,\qquad
 T_\lambda C_\lambda=d_Ud_Vd_W.
\]
Consequently, if all occurrences have indices in $[0,H]$, then
\[
       n_\lambda\le\left\lfloor\frac{H}{T_\lambda}\right\rfloor+1.
\]
\end{theorem}

\begin{proof}
Each $d_Z$ divides a number congruent to one modulo $R$, hence is coprime to
$R$.  Subtracting the two arm values and cancelling $R$ gives
$d_Z\mid |n-m|A_Z$.  Dividing $d_Z$ and $A_Z$ by their gcd shows
$r_Z\mid|n-m|$.  The $r_Z$ remain pairwise coprime because they divide the
pairwise-coprime components $d_Z$, so their product divides the gap.  The
factorization follows from
$(d_Z/\gcd(d_Z,A_Z))\gcd(d_Z,A_Z)=d_Z$.  All indices lie in one residue
class modulo the positive integer $T_\lambda$, which gives the capacity.
\end{proof}

Assume $s>0$ and $N>0$, put $L=\max(s,t)$ and $K=(B+1)(C+1)$, and suppose
$HL\le N$.
Then every direction coefficient is positive and
\[
 C_\lambda\le s(s+Ct)(s+Bt)\le KL^3.
\]

\begin{theorem}[Nonconstant-ray cube-root cap]
\label{thm:affine-ray-cube-root-cap}
If $N^2<d_Ud_Vd_W$, then
\[
                         (n_\lambda-1)^3<KN.
\]
For any finite catalogue $\mathcal L$ of such labels and nonnegative integer
weights $w_\lambda$,
\[
 \sum_{\lambda\in\mathcal L}w_\lambda n_\lambda^3
 \le4(B+1)(C+1)N\sum_{\lambda\in\mathcal L}w_\lambda.
\]
\end{theorem}

\begin{proof}
Put $a=n_\lambda-1$ and $T=T_\lambda$.  The period capacity gives
$aT\le H$, hence $aTL\le N$.  Since
\[
 N^2<TC_\lambda\le TKL^3,
\]
squaring the first inequality and cancelling positive factors gives
$a^2T<KL$.  Thus $a^2<KL$ and $aL\le N$; multiplication yields
$a^3<KN$.  For positive integers $n,X$, the implication
$(n-1)^3<X\Rightarrow n^3\le4X$ follows from
\[
 4(a^3+1)-(a+1)^3=3(a-1)^2(a+1)\ge0.
\]
Apply this with $X=KN$, multiply by $w_\lambda$, and sum.  The constant four
is optimal at this integer level: $n=X=2$ defeats the constant three.
\end{proof}

Ambient positivity is indispensable.  The pure ledger
\[
             (n,H,T,C,D,N,K,L)=(1,0,1,1,1,0,1,1)
\]
satisfies every other period,
factorization, capture, and strict large-label premise, while the proposed
conclusion becomes $0<0$.  This refutes exactly the $N>0$-omitted
strengthening.

Taking
$w_\lambda=\varphi(d_U)\varphi(d_V)\varphi(d_W)$ converts the collinear part
of the third-gcd identity to a total totient-weight problem.  The remaining
cost is the weight of the actual large-label catalogue and the number of
supporting rays, rather than the raw length of a line.

The hypothesis $s>0$ is essential.  On a vertical ray the $U$ component is
constant and may be absorbed in full.  If $d_U\le X_U$, $B,C,t>0$, and
$Ht\le N$, then
\[
 C_\lambda\le X_UBC\,t^2,\qquad
 (n_\lambda-1)^2<X_UBC.
\]
Signed analogues for the two negative-slope arm-level directions, together
with their individual component caps, remain part of the active route.

There is a full actual pressure test.  For the primitive seed $(1,8,9)$ take
\[
 R=6,\quad N=22142,\quad x(n)=(21480,282+n),
 \quad\lambda=(128881,49,121)
\]
and $S=\{0,5929,11858,17787\}$.  All four points lie in the canonical box,
satisfy the admissibility conditions, and carry the displayed pairwise-
coprime label, whose product exceeds $N^2$.  Here
\[
 T_\lambda=5929,\qquad C_\lambda=128881,\qquad
 |S|=17787/5929+1=4,
\]
whereas $\lfloor17787/(d_Ud_Vd_W)\rfloor+1=1$.  This refutes the stronger
claim that the whole label product is itself a ray period and also refutes
deleting the nonconstant-direction premise.  It attains the correct period
capacity and does not contradict the cube-root or vertical theorems.

The formal companion proves all period, capacity, cube-root, weighted-energy,
and vertical statements above and checks this actual counterexample.  No
full-premise counterexample is known to the remaining catalogue-weight and
signed arm-level transfer problems, so those routes remain active.
